\chapter*{Abstrak}
Infrastruktur jalan yang memadai merupakan prasyarat krusial bagi mobilitas ekonomi dan keselamatan transportasi. Inspeksi kerusakan jalan secara manual dinilai tidak efisien, padat karya, dan rentan terhadap subjektivitas. Meskipun algoritma \textit{deep learning} dari keluarga YOLO telah terbukti efektif dalam otomasi deteksi kerusakan jalan, penerapannya di lapangan dihadapkan pada tantangan kesalahan deteksi (\textit{false positive}) akibat objek pengecoh visual (seperti penutup saluran air dan jalan bertambal), serta kebutuhan efisiensi komputasi tinggi pada perangkat terbatas untuk pengawasan waktu nyata (\textit{real-time}).

Penelitian ini melakukan evaluasi komparatif terhadap lima arsitektur deteksi objek berskala nano (`n'): YOLOv8n (\textit{baseline}), arsitektur mutakhir (YOLOv12n dan YOLOv26n), serta dua model rekonstruksi modifikasi domain-spesifik, yaitu YOLOv8-PD (\textit{efficiency-oriented}) dan YOLO-RD (\textit{accuracy-oriented}). Seluruh model dilatih menggunakan dataset N-RDD2024 subset Asia (enam kelas target: empat kerusakan struktural dan dua objek pengecoh) dengan skema pembobotan teradaptasi \textit{Effective Number of Samples} untuk menangani ketimpangan kelas. Kinerja deteksi dievaluasi menggunakan metrik mAP@0.5, mAP@0.5:0.95, dan \textit{macro F1}, sedangkan kelayakan komputasi dianalisis berdasarkan jumlah parameter, GFLOPs, dan kecepatan inferensi (FPS) pada GPU NVIDIA RTX 4070 SUPER. Selain itu, dilakukan studi ablasi pemodelan 6-kelas versus 4-kelas untuk menguji hipotesis penekanan \textit{false positive} menggunakan uji signifikansi statistik \textit{two-proportion z-test}.

Hasil eksperimen membuktikan konfirmasi filosofi desain masing-masing arsitektur. Model rekonstruksi YOLO-RD berhasil meraih akurasi tertinggi dengan nilai test mAP@0.5 sebesar 0,661 dan test mAP@0.5:0.95 sebesar 0,371, membuktikan keunggulan integrasi atensi deformabel dan dekomposisi Haar wavelet. Sementara itu, YOLOv8-PD terbukti paling efisien dengan bobot teringan (2,46M parameter) dan kecepatan inferensi tertinggi mencapai 368,4 FPS (+67,7\% lebih cepat dari baseline) dengan akurasi setara baseline (0,646 vs. 0,648). Uji statistik ablasi pengecoh menunjukkan perbedaan laju \textit{false positive} yang belum signifikan ($p > 0,05$), di mana penambahan kelas pengecoh memperketat batas keputusan visual yang sedikit menekan mAP lubang jalan (\textit{pothole}) namun meningkatkan keseimbangan skor F1 secara agregat. Temuan ini memberikan panduan empiris yang komprehensif dalam pemilihan arsitektur deteksi kerusakan jalan untuk skenario implementasi praktis.

\vspace{1em}
\textbf{Kata Kunci:} Deteksi Kerusakan Jalan, YOLO, Kelas Pengecoh, Rekonstruksi Arsitektur, N-RDD2024, \textit{Lightweight Model}, Efisiensi Komputasi.
